% RESULTADO_RECUPERADO desde II REV09. Véase MAPA_PROCEDENCIA.json.
% Selección de líneas y cambios mecánicos explícitos; ninguna prueba nueva.
\section{Incidence hexadique--octadique et fonctionnelle de Barbero--Immirzi}
\label{v:ant:sec:barbero}
Soit $H\hookrightarrow O$ l’inclusion marquée d’une hexade dans une octade,
$|H|=6$, $|O|=8$, construite au moyen du résidu binaire de la
section~\ref{exc:residuo-binario}. On conserve l’hexade, la dodécade
et le système de coordonnées régional qui déterminent cette inclusion.
Les objets d’incidence
\[
\mathcal F_6=H\times\binom H2,\qquad
\mathcal F_8=O\times\binom H2
\]
ont pour cardinaux $90$ et $120$. Leur différence normalisée donne
\[
\frac{|\mathcal F_6|-|\mathcal F_8|}{24\binom62}=-\frac1{12}.
\]
\begin{figure}[htbp]
\centering
\begin{tikzpicture}[>=Latex,every node/.style={font=\small}]
\draw[rounded corners,draw=ink] (-1.8,-1.35) rectangle (1.8,1.35);
\foreach \a in {0,60,...,300}{
\fill[ink] (\a:0.83) circle(2pt);}
\node at (0,-1.7) {$H:\ 6$ puntos};
\draw[->,thick,ink] (2.1,0)--(3.6,0);
\node[above] at (2.85,1.5) {inclusion marquée};
\draw[rounded corners,draw=ink] (4,-1.35) rectangle (7.6,1.35);
\foreach \a in {0,60,...,300}{
\fill[ink] ($(5.8,0)+(\a:0.83)$) circle(2pt);}
\fill[accent] (4.5,0) circle(2pt);
\fill[accent] (7.1,0) circle(2pt);
\node at (5.8,-1.7) {$O:\ 8$ puntos};
\node at (0,-2.35) {$6\binom62=90$};
\node at (5.8,-2.35) {$8\binom62=120$};
\end{tikzpicture}
\caption{Incidence marquée de six à huit points. Le dessin représente
les ensembles et l'inclusion ; il ne prétend pas transformer les faces ou sommets
d'un cube en blocs de Steiner sans leur application d'incidence.}
\end{figure}

\subsection{Les poids précèdent le coefficient du pôle}
Dans le module libre des événements orientés, on fixe la chaîne
fondamentale du tour
\[
c_{12}^{+}=
\sum_{j\in\Z/12\Z}[j,\mathcal F_6]
-[\partial_{\rm ret},\mathcal F_8].
\]
La multiplicité douze correspond aux secteurs et la multiplicité un
au raccordement orienté. L'involution change l'orientation de la chaîne
complète. Soit
\[
S_m(q)=\frac{q^m}{1-q^{3m}},\qquad 0<q<1.
\]
Le lecteur linéaire attribue $S_{90}$ à chaque secteur et $S_{120}$ au raccordement ;
par conséquent,
\[
\mathscr R(c_{12}^{+})=12S_{90}-S_{120}.
\]

\begin{definition}[Évaluation orientée de l'incidence hexadique--octadique]
\label{paper:barbero-definicion}
Le paramètre \(\gamma_{\HMT}\) est l’évaluation conjuguée de la chaîne
\(c_{12}^{+}\), complétée par sa composante angulaire. Dans le domaine
\(0<q_\pm<1\), il est défini par
\[
\gamma_{\HMT}=\frac{\pi AC^*}{180}
+[\mathscr R(c_{12}^{+})](q_-)-[\mathscr R(c_{12}^{+})](q_+).
\]
L'incidence marquée détermine les espaces d'événements ; la chaîne orientée fixe les multiplicités, et les canaux de la section~\ref{iv:angular:inversion} déterminent l'évaluation.
\end{definition}

\begin{theorem}[Coefficient du pôle et évaluation conjuguée]
L'image de la chaîne fondamentale a pour coefficient $1/24$
devant $(1-q)^{-1}$ et pour résidu $-1/24$ devant $(q-1)^{-1}$.
Son évaluation dans les canaux préalablement engendrés détermine
\begin{equation}
\gamma_{\HMT}=
\frac{\pi AC^*}{180}
+12\bigl[S_{90}(q_-)-S_{90}(q_+)\bigr]
-\bigl[S_{120}(q_-)-S_{120}(q_+)\bigr].
\label{v:ant:eq:gamma}
\end{equation}
\end{theorem}
\begin{proof}
La factorisation $1-q^{3m}=(1-q)\sum_{j=0}^{3m-1}q^j$ implique
$\lim_{q\uparrow1}(1-q)S_m(q)=1/(3m)$. Par linéarité,
\[
\frac{12}{270}-\frac1{360}=\frac1{24}.
\]
La coordonnée $q-1$ inverse le signe du résidu.
L’évaluation conjuguée et la composante angulaire de la définition~\ref{paper:barbero-definicion}
produisent \eqref{v:ant:eq:gamma}.
\end{proof}

Le certificat diophantien $4a-3b=45$ admet les couples
$(12+3t,1+4t)$, $t\ge0$ ; il certifie la minimalité de $(12,1)$,
mais ne remplace pas l'origine des multiplicités dans la chaîne.
L'évaluation de \eqref{v:ant:eq:gamma} est
\[
\gamma_{\HMT}
=0.27400767269445927474725526077398041940\ldots.
\]
Ce chiffre est une sortie de la fonctionnelle orientée ; ce n’est pas une valeur externe
utilisée pour sélectionner $A$, $C^*$ ou leurs coefficients.

\subsection{Parité et sensibilité structurelle}
Inverser $C^*$ échange $q_+$ et $q_-$ et change le signe
de la composante angulaire. Ainsi,
\[
\gamma_{\HMT}(A,-C^*)=-\gamma_{\HMT}(A,C^*).
\]
La fonctionnelle conserve l'orientation, tandis que d'autres lecteurs,
comme $q_+q_-$, sont pairs. Cette différence précède son interprétation physique.
